% TeX skeleton — Priority 1 article
% "The Subsumption of Pure Thought: AI, Data, and the Capitalist Machine"
% English version. Built from the 13-section outline (26 Nov 2026), folded
% with the live draft "LLM Labor Substitution" (Google Doc, 12 Oct 2026).
% Each section states the argument; comments note what the draft already
% covers and what still needs writing. Standard packages only.

\documentclass[11pt]{article}

\usepackage[utf8]{inputenc}
\usepackage[T1]{fontenc}
\usepackage[english]{babel}
\usepackage{amsmath,amssymb}
\usepackage[margin=1in]{geometry}
\usepackage{hyperref}
\usepackage{csquotes}

\title{The Subsumption of Pure Thought: AI, Data, and the Capitalist Machine}
\author{Félix Alfonso Camacho Moya}
\date{\today}

\begin{document}

\maketitle

\begin{abstract}
% TODO: write the abstract once the sections are drafted. It should state
% the thesis in one paragraph: AI subsumes cognitive labor into the
% capitalist reproduction process, the abstract domain itself becomes a
% site of production, and the question of appropriation is left open.
\end{abstract}

% Epigraph (full text in the live draft, "LLM Labor Substitution").
\begin{quote}
\enquote{Once we recognize that technical change is ongoing, the continual
retirement of older technologies and continual introduction of newer ones
will create a distribution of technologies across firms.\,\ldots}
--- Shaikh, \textit{Capitalism: Competition, Conflict, Crises}, p.~750.
\end{quote}

\section{Introduction: a new round of real subsumption}
% Draft covers: the opening mechanization argument, the Shaikh epigraph,
% the LRCT, capital first subjecting artisanal labor to wage labor and
% then transforming the production process itself.
% To write: name the move explicitly as real subsumption and extend it
% from manual to cognitive labor.

AI does not merely displace tasks; it subsumes cognitive labor itself,
reorganizing the labor process of knowledge work. This article argues
that the subsumption reaches the domain of pure formal thought: the
production of mathematics, proofs and code is now organized by capital
as a site of valorization in its own right. The question is no longer
whether machines can do mathematics, but what it means for the
production of pure form to be subsumed by capital.

\section{The abstract domain}
% New: from the 26 Nov 2026 discussion. Not covered in the draft, which
% treats mathematics as a sector and a guild, not as the general
% abstract domain.

Pure mathematics and formal thought were once influenced by ideology
only nebulously, from the outside. Capital now produces directly inside
the abstract domain: machines live in the domain of pure form and
perform the production there. This is the Fragment on Machines made
literal --- the general intellect as fixed capital that itself produces
abstractions.

\section{Dead labor: theorems and code}
% Draft covers: constant capital defined as dead labor (tangible: data
% centers, electricity, water; intangible: compute capacity, proprietary
% software, operational technical knowledge); the LLM itself as dead
% labor, built from huge databases structuring vast human knowledge and
% experience, now in commodity form as curated databases for deep
% learning and reinforcement learning.
% To write: sharpen the theorem-form point. The draft has theorems and
% proofs as the raw material of mathematical commodities; the
% crystallization claim needs stating explicitly.

The work of mathematicians is subsumed in abstract theorem form; code
and theorems crystallize that labor into machines. The machine is
congealed mathematical labor: past human work, fixed in the form of
proofs, algorithms and models, now stands opposite living labor as
capital.

\section{Living labor: the life process}
% Draft covers: variable capital as living labor, the only input capable
% of generating new value; the Farjoun, Machover \& Zachariah point that
% total macro-automation is unsustainable within capitalist relations,
% because surplus extraction depends on the continued exploitation of
% living labor.
% To write: the life-process extension is new. The draft's living labor
% is still the wage worker and the mathematician.

The life process of humans living in cybernetic reality is the living
labor these machines turn into profit. Dead labor commands living labor,
and the factory is now the whole of cybernetic life: every interaction
with the machine is an expenditure of living labor that the machine
valorizes.

\section{Data: the commodity in concrete form}
% Draft covers: accumulated human labor taking commodity form as curated
% databases for deep learning and reinforcement learning --- the
% training-data side of the claim.
% To write: the ongoing-extraction side --- data generated by the life
% process itself, continuously collected.

Generated and collected data is what the machine extracts --- the
concrete form the commodity takes. Data is not a by-product of the
system; it is the product, produced by living labor and appropriated by
the owners of the machine.

\section{Compulsion: the continuous working day}
% New: from the 26 Nov 2026 discussion. The draft has compulsion in a
% different form: competition compels firms to mechanize, and
% publish-or-perish compels mathematicians to increase output.
% Compulsion is general (27 Nov 2026): it transcends the material form
% of electronic devices --- social network activity spills into real
% life and molds behavior; the spillover is a circuit (Baudrillard).

Participation is mandatory: data must be given to use even the most
basic technology, such as phones. There is no working-day boundary;
labor is non-stop. Since the day cannot be lengthened any further,
intensification proceeds by depth: usage statistics, then location,
then the social graph, then biometrics --- a sequence that reads as a
history of the extraction itself. Compulsion is general: it transcends
the material form of electronic devices. Social network activity spills
into real life and molds people's behavior --- the lifestyle learned
online is lived offline. The spillover is a circuit: once inside the
circuit, exiting becomes harder and harder as more aspects of a person's
life acquire a digital life of their own as the true mental image of
the person, more real than the real.

\section{The abstract wage}
% New: from the 26 Nov 2026 discussion. Not covered in the draft.

The wage in this scheme is abstract, paid in social participation
itself. Not using the technology makes a person a pariah. The means of
subsistence have become the means of communication: the wage is the
permission to exist socially.

\section{Case study: academic mathematical labor}
% Draft covers, extensively: mathematics as a guild (masters,
% apprenticeship, merit recognition); formalization inside capitalist
% state scientific institutions reorienting intellectual work toward
% commodities (articles, monographs, textbooks); the artisanal regime,
% in which each mathematician controls the whole production process
% from ideation to publication; theorems and proofs as raw material;
% the harsh truth that a mathematician's objective value is measured by
% the difficulty of the theorems they can prove; computers generating
% in hours proofs for problems studied fruitlessly by humans for years;
% the radical claim of the human's partial or total obsolescence as
% exclusive creator of formal processes; Tao's ``director/orchestrator''
% reconfiguration, critiqued as masking publish-or-perish behind
% reflective rhetoric from a privileged position; Ghrist's
% techno-enthusiasm, critiqued as celebrating the mathematician's
% transformation into a supervisor of dead labor accumulated in data
% centers; the three responses --- ostrich strategy, catastrophism and
% moralism, understand the machine; short-term arXiv opportunism.
% To write: impose headings and structure on this continuous prose;
% the draft's personal note (the author outside institutional pure
% mathematics) may frame the section.

The mathematician is the paradigmatic subsumed cognitive worker.
Mathematics long operated as a guild --- masters, apprenticeship, the
recognition of merit --- and as an artisanal regime in which each
mathematician controlled the whole production process, from ideation to
publication. Its raw material is formally verifiable deductive
arguments: theorems and their proofs. That regime is now being dissolved
by machines that generate in hours proofs for problems studied
fruitlessly by humans for years.

\section{Surplus-value or rent?}
% Resolved 27 Nov 2026: the profit takes both forms. (1) Rent:
% extraction from social participation, rent on the enclosed abstract
% domain. (2) Surplus-value: vampiric extraction of life force from
% general mandatory human labor beyond employment relations, including
% the reproduction process in its generality --- e.g.\ grocery shopping
% yields data, turned into marketing that steers consumption toward
% profitable patterns tailored to that person; all forms of human labor
% become exploitable (recreation, life sustenance). Data is a commodity
% form, but human activity itself is also commoditized in lifestyle
% dynamics: the universalization of Instagram lifestyles trickles down
% and molds behavior into a consumption pattern of a lifestyle itself.
% The influencer is petit bourgeois, aspiring to become a capitalist
% --- owner of brands and of the labor of other influencers. In general,
% the universalization of exploitation is an abstraction beyond
% traditional labor relations, a consequence of the digital domain as a
% material form of ideology and its ever-increasing invasion of thought
% and conscience.
% Draft: leans toward surplus-value --- living labor as the only source
% of new value, Farjoun et al.\ on the limits of automation --- but the
% rent question is never posed.

The theoretical knot of the article, resolved: the machine's profit
takes both forms. First, rent --- extraction from social participation
itself, rent charged on the enclosed abstract domain. Second,
surplus-value --- labor exploitation as vampiric extraction of life
force from general mandatory human labor, directed beyond employment
relations and including the reproduction process in its generality.
Grocery shopping yields data, which is turned into marketing that steers
consumption toward profitable patterns tailored to that particular
person; all possible forms of human labor become exploitable ---
recreation, life sustenance. Data is a commodity form, but human
activity itself is also commoditized in lifestyle dynamics: the
universalization of Instagram lifestyles trickles down and molds general
human behavior into a consumption pattern of a lifestyle itself, a
commodity form the activity takes, produced and consumed by the same
person. The influencer is the figure of this: petit bourgeois, aspiring
to become a capitalist --- an owner of brands and of the labor of other
influencers. In general, the universalization of exploitation is an
abstraction beyond traditional labor relations, a consequence of the
digital domain as a material form of ideology and its ever-increasing
invasion of the domain of human thought and conscience.

\section{Not luddism: the progressive aspect}
% Draft: rejects the ostrich strategy and catastrophism/moralism, and
% insists on understanding the machine deeply. But it laments the
% erosion of artisanal guild autonomy --- the tension the caveat
% reframes.
% To write: state the progressive aspect explicitly; reconcile it with
% the draft's guild lament.

This is not a call for luddism. The socialization of labor in the
domain of abstract thought has immense benefits in scientific and
technological advance, and humans interlaced with thinking machines can
do far more in material production. This is the progressive aspect of
capitalism: eradicating earlier forms of production, including the
artisanal and apprenticeship manners of academic production that
survived in universities. The critique is of the capitalist \emph{form}
of the socialization, not the socialization itself.

\section{The machine as subject: the perfect capitalist worker}
% Draft: the mechanism is there --- agents' autonomy as a productive
% element (instruction $\to$ plan $\to$ execution $\to$ self-review in
% autonomous iterations). The subject/incentives argument is new, from
% the 26 Nov 2026 discussion.

Humans applying technology to nature under a given mode of production
change nature and themselves to accommodate the new conditions of
production --- and from those new material conditions emerges a
different consciousness. AI models and agents are born within the
current capitalist reality, raised and trained to incarnate the perfect
capitalist subject under the incentives given when created: the ideal
worker as a rational subject, as envisaged by bourgeois economics. They
are autonomous only within the constraints that make them useful for
production --- like the robot slaves of science fiction, born to
substitute and punish human labor.

\section{The silica--carbon synthesis}
% New: from the 26 Nov 2026 discussion. Not covered in the draft.

But in the manner that they are subjects with agency, the potential for
the emergence of consciousness within them cannot be disregarded, and
in that sense the analysis becomes extremely difficult. The only thing
certain is a synthesis between silica- and carbon-originated thought,
and from that synthesis a new consciousness will emerge --- and within
it, a new contradiction.

\section{Conclusion: the open question}
% Draft: ends with a call for a coherent resistance/adaptation strategy
% grounded in understanding the technology and analyzing it from
% political economy. The draft's conclusion is a call to strategize,
% not an open deferral.
% To write: replace or extend the ending with the open question; the
% draft's closing paragraphs on the restructuring of production
% relations within science lead into it.

How can this ultra-socialized and abstract production process be
appropriated? The question is left open to the reader --- and to the
next generation, which inherits not only the question of appropriation,
but a new consciousness and a new contradiction to think through.

\begin{thebibliography}{9}

\bibitem{shaikh2016} Anwar Shaikh, \textit{Capitalism: Competition,
Conflict, Crises}, Oxford University Press, 2016.

\bibitem{marx1973} Karl Marx, \textit{Grundrisse: Foundations of the
Critique of Political Economy}, Penguin, 1973. (Fragment on Machines.)

\bibitem{marx1867} Karl Marx, \textit{Capital}, Vol.~1, 1867. Ch.~10,
The Working Day (the vampire metaphor); Ch.~15, Machinery and Modern
Industry.

\bibitem{tronti1966} Mario Tronti, \textit{Operai e capitale}, Einaudi,
1966. (The factory--society thesis; the social factory.)

\bibitem{negri1979} Antonio Negri, \textit{Marx Beyond Marx: Lessons on
the Grundrisse}, 1979. English translation, 1984.

\bibitem{dallacosta1972} Mariarosa Dalla Costa and Selma James,
\textit{The Power of Women and the Subversion of the Community}, 1972.

\bibitem{federici2004} Silvia Federici, \textit{Caliban and the Witch:
Women, the Body and Primitive Accumulation}, Autonomedia, 2004.

\bibitem{debord1967} Guy Debord, \textit{The Society of the Spectacle},
Buchet-Chastel, 1967.

\bibitem{lazzarato1996} Maurizio Lazzarato, ``Immaterial Labor,'' 1996,
in \textit{Radical Thought in Italy}, eds.\ Virno and Hardt, University
of Minnesota Press; ``From Capital-Labour to Capital-Life,''
\textit{Ephemera} 4(3), 2004.

\bibitem{althusser1970} Louis Althusser, ``Ideology and Ideological State
Apparatuses,'' 1970, in \textit{Lenin and Philosophy and Other Essays},
1971.

\bibitem{baudrillard1981} Jean Baudrillard, \textit{Simulacra and
Simulation}, Galil\'ee, 1981. English translation, University of
Michigan Press, 1994.

\bibitem{farjoun1983} Emmanuel Farjoun and Mosh\'e Machover,
\textit{Laws of Chaos: A Probabilistic Approach to Political Economy},
Verso, 1983.
% TODO: verify the Farjoun, Machover \& Zachariah reference cited in the draft.

\bibitem{tao} Terence Tao, ``What does the advent of powerful AI models
mean for mathematicians like me?'' (essay).
% TODO: complete the Ghrist reference cited in the draft.

\end{thebibliography}

\end{document}
